%==============================================================================
% KAPITEL 0: ORIENTIERUNG
% Einzufügen nach \tableofcontents und vor \section{Einführung}
%==============================================================================

\section*{Orientierung: Über dieses Dokument}
\addcontentsline{toc}{section}{Orientierung: Über dieses Dokument}

%------------------------------------------------------------------------------
\subsection*{Executive Summary}
\addcontentsline{toc}{subsection}{Executive Summary}
%------------------------------------------------------------------------------

\begin{tcolorbox}[colback=blue!5!white, colframe=blue!75!black, title=Das Wichtigste in 60 Sekunden]

\textbf{Was ist das Problem?}\\
Die Verhaltensökonomie hat in 50 Jahren hunderte Effekte dokumentiert (Loss Aversion, Framing, Anchoring...), aber nie ein kohärentes Fundament gebaut. Das macht systematische Anwendung unmöglich.

\textbf{Was ist die Lösung?}\\
Das \textbf{Behavioral Change Model (BCM)} formalisiert menschliches Verhalten in einem axiomatischen Framework mit 115 expliziten Annahmen. Das \textbf{META-Modul} dokumentiert diese Fundamente.

\textbf{Was ist neu?}\\
\textbf{BEATRIX} nutzt Large Language Models (LLMs), um BCM für Nicht-Experten zugänglich zu machen: Natürliche Sprache rein, quantifizierte Verhaltensprognosen raus.

\textbf{Warum ist das wichtig?}\\
Erstmals können verhaltensökonomische Analysen \textit{skalierbar}, \textit{konsistent} und \textit{nachvollziehbar} durchgeführt werden -- von der Persona-Beschreibung bis zur Handlungsempfehlung.

\end{tcolorbox}

%------------------------------------------------------------------------------
\subsection*{Was ist BCM/BEATRIX?}
\addcontentsline{toc}{subsection}{Was ist BCM/BEATRIX?}
%------------------------------------------------------------------------------

\paragraph{Der Elevator Pitch}

\begin{quote}
\textit{``BCM ist für menschliches Verhalten, was die Finite-Elemente-Methode für Ingenieure ist: Ein mathematisches Framework, das komplexe Realität in berechenbare Modelle übersetzt. BEATRIX ist die Benutzeroberfläche, die dieses Framework für Nicht-Mathematiker zugänglich macht.''}
\end{quote}

\paragraph{Ein konkretes Beispiel}

\textbf{Frage:} ``Wie reagiert Anna (45, Lehrerin, umweltbewusst, preissensitiv) auf eine 5\% Preiserhöhung unseres Produkts?''

\textbf{Ohne BCM/BEATRIX:}
\begin{itemize}
    \item Befragung (teuer, langsam, Bias-anfällig)
    \item Bauchgefühl des Produktmanagers (HiPPO-Problem)
    \item Historische Daten (falls vorhanden)
\end{itemize}

\textbf{Mit BEATRIX:}
\begin{enumerate}
    \item LLM extrahiert BCM-Parameter aus Persona-Beschreibung ($\lambda = 2.3$, $w_{\text{eco}} = 0.4$, ...)
    \item BCM berechnet Churn-Wahrscheinlichkeit: $4.2\%$ [CI: 2.8\%--5.9\%]
    \item LLM erklärt: ``Anna's hohe Verlustaversion ($\lambda = 2.3$) macht sie preissensitiver als der Durchschnitt, aber ihre ökologischen Werte könnten bei `grüner' Kommunikation als Puffer wirken.''
\end{enumerate}

\paragraph{Was BCM/BEATRIX NICHT ist}

\begin{itemize}
    \item \textbf{Keine Kristallkugel:} BCM liefert probabilistische Prognosen mit Konfidenzintervallen, keine Gewissheiten
    \item \textbf{Kein Ersatz für Empirie:} BCM-Prognosen sollten durch Feldexperimente validiert werden
    \item \textbf{Kein Manipulationswerkzeug:} Die Governance-Axiome (Kapitel 4) definieren ethische Grenzen
    \item \textbf{Keine Black Box:} Jede Prognose ist durch Axiome und Parameter nachvollziehbar
\end{itemize}

%------------------------------------------------------------------------------
\subsection*{Für wen ist dieses Dokument?}
\addcontentsline{toc}{subsection}{Für wen ist dieses Dokument?}
%------------------------------------------------------------------------------

\begin{table}[h]
\centering
\caption{Zielgruppen und Leseempfehlungen}
\begin{tabular}{@{}p{3.5cm}p{5cm}p{3cm}p{2cm}@{}}
\toprule
\textbf{Zielgruppe} & \textbf{Interesse} & \textbf{Kapitel} & \textbf{Lesezeit} \\
\midrule
\textbf{Verhaltens-ökonomen} & Theoretische Fundierung, Axiomatik & 1--5, 10 & 2--3 Std \\
\textbf{KI/ML Engineers} & Technische Implementation, LLM-Integration & 6, 7, 9 & 1--2 Std \\
\textbf{Business Stakeholder} & Value Proposition, Anwendungsfälle & 0, 1.9, 10 & 30 Min \\
\textbf{Regulatoren, Ethiker} & Governance, Transparenz, Grenzen & 4, 1.9 & 1 Std \\
\textbf{Berater, Praktiker} & Praxisanwendung, Methodologie & 1.9, 9 & 1 Std \\
\bottomrule
\end{tabular}
\end{table}

%------------------------------------------------------------------------------
\subsection*{Wie lese ich dieses Dokument?}
\addcontentsline{toc}{subsection}{Wie lese ich dieses Dokument?}
%------------------------------------------------------------------------------

\paragraph{Track A: Schnelldurchlauf (30 Minuten)}
Für Entscheider, die verstehen wollen, worum es geht:
\begin{enumerate}
    \item Diese Orientierung (Kapitel 0)
    \item Motivation und Value Proposition (Kapitel 1.1, 1.9)
    \item Zusammenfassung (Kapitel 10)
\end{enumerate}

\paragraph{Track B: Konzeptuelles Verständnis (2--3 Stunden)}
Für Fachleute, die das Framework verstehen wollen:
\begin{enumerate}
    \item Track A, plus:
    \item Die vier Axiom-Kategorien (Kapitel 2--5)
    \item Evidenzhierarchie und Biases (Kapitel 3.5, 3.4)
\end{enumerate}

\paragraph{Track C: Technische Tiefe (4--5 Stunden)}
Für Entwickler und Implementierer:
\begin{enumerate}
    \item Track B, plus:
    \item Integration mit operativen Modulen (Kapitel 6)
    \item META im BEATRIX-System (Kapitel 7)
    \item Von Handwerk zu Engineering (Kapitel 9)
\end{enumerate}

\paragraph{Track D: Vollständig (1 Tag)}
Für Wissenschaftler und tiefe Analyse:
\begin{enumerate}
    \item Alle Kapitel sequentiell
    \item Vollständige Axiom-Listen im Anhang
    \item Literaturverzeichnis
\end{enumerate}

%------------------------------------------------------------------------------
\subsection*{Dokumentstruktur auf einen Blick}
\addcontentsline{toc}{subsection}{Dokumentstruktur auf einen Blick}
%------------------------------------------------------------------------------

\begin{figure}[h]
\centering
\begin{tikzpicture}[
    node distance=0.8cm,
    chapter/.style={rectangle, draw, rounded corners, minimum width=4cm, minimum height=0.8cm, align=center, font=\small},
    arrow/.style={->, thick, >=stealth}
]

% Foundation
\node[chapter, fill=gray!20] (c0) at (0, 0) {0. Orientierung};

% Core Theory
\node[chapter, fill=blue!15] (c1) at (0, -1.2) {1. Einführung \& Motivation};
\node[chapter, fill=green!15] (c2) at (-4, -2.4) {2. Ontologie};
\node[chapter, fill=green!15] (c3) at (-1.3, -2.4) {3. Epistemologie};
\node[chapter, fill=green!15] (c4) at (1.3, -2.4) {4. Governance};
\node[chapter, fill=green!15] (c5) at (4, -2.4) {5. Dynamik};

% Integration
\node[chapter, fill=orange!15] (c6) at (-2, -3.6) {6. Integration};
\node[chapter, fill=orange!15] (c7) at (2, -3.6) {7. META in BEATRIX};

% Transformation
\node[chapter, fill=purple!15] (c9) at (0, -4.8) {9. Transformation};

% Conclusion
\node[chapter, fill=gray!20] (c10) at (0, -6) {10. Zusammenfassung};

% Arrows
\draw[arrow] (c0) -- (c1);
\draw[arrow] (c1) -- (c2);
\draw[arrow] (c1) -- (c3);
\draw[arrow] (c1) -- (c4);
\draw[arrow] (c1) -- (c5);
\draw[arrow] (c2) -- (c6);
\draw[arrow] (c3) -- (c6);
\draw[arrow] (c4) -- (c7);
\draw[arrow] (c5) -- (c6);
\draw[arrow] (c6) -- (c9);
\draw[arrow] (c7) -- (c9);
\draw[arrow] (c9) -- (c10);

% Legend
\node[right=0.5cm of c5, font=\footnotesize, align=left] {
    \textcolor{gray!60}{$\blacksquare$} Rahmen\\
    \textcolor{blue!50}{$\blacksquare$} Motivation\\
    \textcolor{green!50}{$\blacksquare$} Axiom-Kategorien\\
    \textcolor{orange!50}{$\blacksquare$} Integration\\
    \textcolor{purple!50}{$\blacksquare$} Transformation
};

\end{tikzpicture}
\caption{Kapitelstruktur und Abhängigkeiten des META-Modul Papers}
\end{figure}

%------------------------------------------------------------------------------
\subsection*{Konventionen und Notation}
\addcontentsline{toc}{subsection}{Konventionen und Notation}
%------------------------------------------------------------------------------

\paragraph{Axiom-Nummerierung}

Alle Axiome folgen dem Schema \texttt{MA-XXX-\#\#}:
\begin{itemize}
    \item \texttt{MA} = Meta-Axiom
    \item \texttt{XXX} = Kategorie (ONT, EPI, GOV, DYN, INF)
    \item \texttt{\#\#} = Laufende Nummer
\end{itemize}
Beispiel: \texttt{MA-EPI-12} = Meta-Axiom, Kategorie Epistemologie, Nummer 12 (Availability Bias)

\paragraph{Abkürzungen}

\begin{table}[h]
\centering
\small
\begin{tabular}{@{}ll|ll@{}}
\toprule
\textbf{Kürzel} & \textbf{Bedeutung} & \textbf{Kürzel} & \textbf{Bedeutung} \\
\midrule
BCM & Behavioral Change Model & META & Meta-Axiome Modul \\
INU & Individual Utility & KNU & Contextual Utility \\
IDN & Identity Utility & $\lambda$ & Loss Aversion Parameter \\
BEATRIX & BCM + LLM System & LLM & Large Language Model \\
FEPSDE & Financial, Emotional, Physical, Social, Digital, Ecological \\
\bottomrule
\end{tabular}
\end{table}

\paragraph{Hervorhebungen}

\begin{itemize}
    \item \textbf{Fettdruck}: Wichtige Begriffe bei Erstnennung
    \item \textit{Kursiv}: Zitate, fremdsprachliche Begriffe
    \item \texttt{Monospace}: Code, Parameter, Axiom-IDs
    \item Farbige Boxen: Kernaussagen, Warnungen, Zusammenfassungen
\end{itemize}

\vspace{1em}
\begin{tcolorbox}[colback=green!5!white, colframe=green!75!black, title=Bereit?]
Sie haben nun alle Informationen, um dieses Dokument effektiv zu nutzen. Wählen Sie Ihren Track (A--D) und beginnen Sie mit Kapitel 1 -- oder springen Sie direkt zu den für Sie relevanten Abschnitten.
\end{tcolorbox}

\newpage
